\scanpage{107}
\mspage{53}

qui se précise par l'hom \unsure{effectif|fonctel} de groupes profinis
\[
(6)\qquad \mathcal E_X^{\widetilde X}\longrightarrow\mathcal E_Y^{\widetilde Y}
\]
où $\widetilde Y$ est le rev. universel de $Y$ défini par $\widetilde X\longrightarrow Y$ ($\widetilde X$ jouant le rôle du point --- on devrait pouvoir écrire en fait $\mathcal E_Y^{\widetilde X}$, plus gén. $\mathcal E_Y^Z$ chaque fois qu'on a un $Y$-schéma $Z$ $1$-connexe, jouant le rôle de ``foncteur fibre'' pour le topos $B_{\unsure{\pi(X)}}$ des rev étales de $Y$\unsure{.})

\deleted{\unsure{que} $\mathcal E_X^{\widetilde X}$}

On peut dire que
\[
(7)\qquad (X,\widetilde X)\longmapsto\mathcal E_X^{\widetilde X}
\]
est un foncteur, de la catégorie des schémas $0$-connexes $X$ \uline{munis} d'un rev. universel (ou seulement d'un $Z$ $1$-connexe s'envoyant dans $X$) vers celle des groupes profinis. Ceci s'applique en particulier en regardant la \inserted{sous-}cat. des $(X,\unsure{\xi|\zeta})$ munis d'un $p^t$ géom ---

\margin{106}
